\section{问题一：忽略库存的推荐算法}\label{sec:q1}

\subsection{算法 1 在有限库存下的两类缺陷}
算法 1 在营业开始前为每类顾客计算一次排序，此后对同类顾客始终推荐相同的 $K$ 件商品。库存有限时，这一做法存在两类缺陷。

第一类是可行性缺陷。商品售罄后，算法 1 仍然推荐它，推荐位被一件无法成交的商品占据：首选该商品的顾客空手离开，平台本可以在这个位置上展示其他有货商品。这一行为同时违反约束\eqref{eq:feasible}，也就是题目评分中“零缺货推荐”的要求。

第二类是稀缺性缺陷。即使先剔除缺货商品，按 $v_{ij}$ 排序仍然只考虑本期收益。一件稀缺商品卖给当前顾客，后续更需要它、又没有替代选择的顾客就失去了购买机会。第一类缺陷通过库存过滤即可修正，得到库存感知贪心策略 \Greedy{}：每期在有货商品中按 $v_{ij}$ 选取前 $K$ 件；第二类缺陷需要估计库存的机会成本，这是问题三的任务。

据此回答问题一：在有限库存下，算法 1 会推荐缺货商品，平台需要先对它做库存过滤才能满足可行性要求，进一步提高收入还需要引入库存的机会成本。下面的算例说明，缺货推荐造成的收入损失可以任意接近全部收入。

\subsection{收入损失超过 30\% 的算例}
\begin{example}\label{ex:q1}
取 $m=1$，$n=3$，$K=2$，$r=(10,8,6)$，$p=(0.5,0.3,0.2)$，$C=(1,1,T)$。
\end{example}

算法 1 依据 $v=(5,\,2.4,\,1.2)$ 每期推荐 $\{1,2\}$。商品 1 与商品 2 各只有一件，售出之后对应的推荐位闲置，商品 3 始终得不到展示。商品 1 在 $T$ 期内被首选至少一次的概率为 $1-0.5^T$，商品 2 为 $1-0.7^T$，于是
\begin{equation}
\E R^{\AlgOne}(T)=10\bigl(1-0.5^{T}\bigr)+8\bigl(1-0.7^{T}\bigr)<18.
\label{eq:q1naive}
\end{equation}
作为对照，考虑一个容易手算的可行策略 $\pi'$：每期推荐商品 3，商品 1 有货时同时推荐商品 1。商品 3 的库存等于期数，永远不会售罄，因此
\begin{equation}
\E R^{\pi'}(T)=10\bigl(1-0.5^{T}\bigr)+6\times0.2\times T.
\label{eq:q1feasible}
\end{equation}

\begin{proposition}\label{prop:q1}
在算例\ref{ex:q1}中，算法 1 相对 $\pi'$ 的收入损失 $L(T)=1-\E R^{\AlgOne}(T)/\E R^{\pi'}(T)$ 满足 $L(T)>1-15/T$，因而 $\lim_{T\to\infty}L(T)=1$；当且仅当 $T\ge14$ 时 $L(T)>30\%$，其中 $L(20)=47.08\%$。
\end{proposition}
\begin{proof}
由式\eqref{eq:q1naive}与式\eqref{eq:q1feasible}，$\E R^{\AlgOne}<18$，$\E R^{\pi'}>1.2T$，故 $L(T)>1-18/(1.2T)=1-15/T$，极限结论随之成立。当 $T\ge22$ 时 $1-15/T>0.3$；对 $T\le21$ 逐一计算可得 $L(T)$ 随 $T$ 递增，$L(13)=29.99\%$，$L(14)=33.04\%$，$L(20)=47.08\%$。
\end{proof}

$\pi'$ 只用于说明损失的量级。表\ref{tab:q1}进一步列出库存感知贪心 \Greedy{}、本文第\ref{sec:policy}节的两种 LP 策略以及最优动态规划的精确期望收入。记最优期望收入为 $V^{*}(T)$，算法 1 相对最优策略的损失为 $L^{*}(T)=1-\E R^{\AlgOne}(T)/V^{*}(T)$，推荐位上出现缺货商品的期望次数为
\begin{equation}
W(T)=\E\sum_{t=1}^{T}\bigl|\{j\in S_t: c_{tj}=0\}\bigr|.
\label{eq:waste}
\end{equation}

\begin{table}[htbp]
\centering\small
\caption{算例\ref{ex:q1}的精确期望收入、算法 1 的相对损失 $L^{*}(T)$ 与缺货推荐次数 $W(T)$}\label{tab:q1}
\begin{tabular}{rrrrrrrrr}
\toprule
$T$ & $\E R^{\AlgOne}$ & $\E R^{\pi'}$ & $\E R^{\Greedy}$ & $\E R^{\LPFD}$ & $\E R^{\LPDual}$ & $V^{*}$ & $L^{*}(T)$ & $W(T)$ \\
\midrule
5 & 16.343 & 15.688 & 20.843 & 20.983 & 20.843 & 21.018 & 22.2\% & 5.3 \\
10 & 17.764 & 21.990 & 28.264 & 29.446 & 29.420 & 29.447 & 39.7\% & 14.8 \\
20 & 17.994 & 34.000 & 40.494 & 41.984 & 41.984 & 41.984 & 57.1\% & 34.7 \\
50 & 18.000 & 70.000 & 76.500 & 78.000 & 78.000 & 78.000 & 76.9\% & 94.7 \\
100 & 18.000 & 130.000 & 136.500 & 138.000 & 138.000 & 138.000 & 87.0\% & 194.7 \\
200 & 18.000 & 250.000 & 256.500 & 258.000 & 258.000 & 258.000 & 93.0\% & 394.7 \\
\bottomrule
\end{tabular}
\end{table}


表\ref{tab:q1}显示：$T=20$ 时，算法 1 平均有 \QoneWasteTwenty{} 个推荐位展示缺货商品（共 40 个推荐位），相对最优策略的损失 $L^{*}(20)=\QoneLossOptTwenty\%$；$T=100$ 时 $L^{*}$ 升至 \QoneLossOptHundred\%。库存过滤消除了绝大部分损失；$T\ge20$ 时 \Greedy{} 与最优值的差距稳定在 1.5，这部分损失来自开局阶段同时推荐商品 1、2 而推迟展示商品 3：商品 3 的库存充裕、机会成本为零，最优策略从第一期起就把它放进推荐列表。两种 LP 策略与最优值的差距在各个 $T$ 下都不超过 \QoneLPMaxGap{}。

\subsection{课程实例上的算法 1}\label{sec:q1course}
课程实例的库存只是轻度紧张，算法 1 的损失相应较小，但缺货推荐依然存在。商品 1 的 Top-2 期望需求为 42 件，高于库存 40 件，多数样本路径在营业后段售罄商品 1；此后类型 1、2 顾客的推荐列表中始终保留商品 1。精确计算表明，五份数据上算法 1 每条路径平均出现 \AlgOneOOS{} 次缺货推荐，期望收入为 \AlgOneER{}，比 \Greedy{} 低 \AlgOneLossGreedy{}。第\ref{sec:sens}节把库存同比例缩小后，算法 1 的缺货推荐次数和收入损失都迅速上升。
